\documentclass[../../main.tex]{subfiles}
\begin{document}

\begin{flushright}
\footnotesize\textit{【自動文字起こし・要確認】}
\end{flushright}

[解] $P(\alpha, \alpha^2), Q(\beta, \beta^2) \; (\alpha < \beta)$ とおく。題意から $PQ$ の方程式 $y = l(x)$ として
\begin{align}
1 &= \int_{\alpha}^{\beta} (l(x) - x^2) dx \\
&= \int_{\alpha}^{\beta} -(x - \alpha)(x - \beta) dx \\
&= \dfrac{1}{6} (\beta - \alpha)^3
\end{align}
$\therefore \beta - \alpha = \sqrt[3]{6} \cdots \text{①}$

このもとで、$R(X, Y)$ の軌跡をもとめる。
$$ X = \dfrac{\alpha + \beta}{2} , \quad Y = \dfrac{\alpha^2 + \beta^2}{2} \cdots \text{②} $$
であり、$t = \alpha + \beta, s = \beta - \alpha$ とおくと、$\alpha \beta = \dfrac{1}{4} (t^2 - s^2)$。①で、②に代入して
$$ X = \dfrac{1}{2} t , \quad Y = \dfrac{t^2 + s^2}{4} = \dfrac{1}{4} \left( t^2 + 6^{\dfrac{2}{3}} \right) \cdots \text{③} $$
又、$\alpha, \beta$ は $x$ の2次方程式 $x^2 - tx + \dfrac{1}{4}(t^2 - 6^{\dfrac{2}{3}}) = 0$ の2実数解。この判別式 $D$ として $D = 6^{\dfrac{2}{3}} > 0$ だから、$\alpha, \beta$ は常に存在する。したがって③より $t$ を消去して、
$$ Y = X^2 + \dfrac{1}{4} 6^{\dfrac{2}{3}} $$

\end{document}
